\scanpage{47}\mspage{46}
% slanted
\margin{Ainsi, $M$ s'identifie à un ens de parties de $\mathcal F$ (savoir les $F_X$) i.e. à l'ens des figures élémentaires. \deleted{Donc on peut considérer que la structure est \unsure{portée} \unsure{par|sur} $\mathcal F$, muni de $\mathcal F_{\unsure{\mathrm{él}|\mathrm{el}}}$ et de $\unsure{\Sigma}$, la relation de subdivision (relation d'ordre)} \deleted{Il y a une sorte de symétrie -- on travaille sur $\mathcal F$, \unsure{en tant|ou bien|comme} \inserted{\unsure{axiomatique des|par les axiomes des}} \unsure{et tantôt} sur $ M$,  1°) relation d'ordre $\unsure{C'|\leq}$ (subdivision, ou $\ll$ raffinement) et ens. de parties (élémentaires, les figures).}}
\deleted{Cela étant}

\deleted{$\unsure{X}\unsure{\ll|\leq}Y\quad\unsure{F'}\unsure{\supset|\subset|\in}F_X$} Si $F$ est une figure, alors $X\in F_{\deleted{\unsure{X|x|Y}}}$ équiv. à : $F_X\subset F$, qui implique que pour tout $Y\in F_X$ \inserted{i.e. $Y\leq X$} on a $F_Y\subset F$ i.e. $Y\in F$. Ainsi, $F$ est une partie de $ M$ \uline{fermée} pour $\leq$.

On dit que deux figures \inserted{$F,G$} sont \uline{compatibles} si $F\cup G$ est une figure, et que $X,Y\in M$ sont \uline{compatibles} si $F_X,F_Y$ le sont

\uline{C \deleted{\unsure{2|3}}\unsure{3|2}} Pour que $F$ et $G$ soient compatibles, il f. et s. que $\forall\,\unsure{x|X}\in F,\ Y\in G$, $X$ et $Y$ le \deleted{\unsure{soi}} \unsure{sont}.
\margin{\tikz{\draw (0,0)--(0,.8);\draw (.06,0)--(.06,.8);}}
% slanted
\margin{correspond \unsure{à|:} Fig 4}

\uline{Cor} \deleted{En particulier si $X$} Pour que $X,Y\in M$ soient compatibles, il f et s. qu'ils a\uuline{pp}artiennent à une figure

\deleted{\uline{C 3} Si $F$ est.}
% slanted
\margin{intervertir \tikz[baseline=(a.base)]{\node[inner sep=1pt] (a) {}; \draw[->] (a.north) .. controls +(0,.6) and +(-.5,0) .. (1.4,1); \draw[->] (a.south) .. controls +(0,-.6) and +(-.5,0) .. (1.4,-1);}}

\uline{C \deleted{\unsure{3|2}}\unsure{2|3}} Soit $F$ une figure, alors tte partie de $F$ fermée pour $\leq$ est \unsure{aussi|ainsi} une figure (i.e. pour qu'une partie $F'$ de $F$ soit une figure, il f et il \uline{suffit} qu'elle soit fermée pour $\leq$).
\margin{\tikz{\draw (0,0)--(0,.8);\draw (.06,0)--(.06,.8);}}

(\uline{Sous-figures}) Si $ M\ne\phi$ cela \deleted{\unsure{signifie}} implique que la partie vide de $ M$ est une figure. Si $ M=\phi$, \deleted{cela implique que} \inserted{on aura $\mathcal F=\phi$} \deleted{si $\mathcal F$} \unsure{ou|on} $\mathcal F=\{\phi\}$, et on suppose que c'est le deuxième cas : \uline{figure vide} existe tjrs
% slanted
\margin{\uline{Contrée vide} si $ M=\phi$}
